% Generated draft: requires mathematical and source review.
\documentclass[11pt]{article}
\usepackage[margin=25mm]{geometry}
\usepackage{fontspec}
\setmainfont{Latin Modern Roman}
\usepackage{amsmath,amssymb,mathtools,unicode-math}
\setmathfont{Latin Modern Math}
\usepackage{enumitem,needspace,xurl,hyperref}
\setlength{\parindent}{0pt}
\setlength{\parskip}{5pt}
\setlength{\emergencystretch}{3em}
\newcommand{\N}{\mathbb{N}}
\newcommand{\Z}{\mathbb{Z}}
\newcommand{\Q}{\mathbb{Q}}
\newcommand{\R}{\mathbb{R}}
\newcommand{\C}{\mathbb{C}}
\newcommand{\F}{\mathbb{F}}
\newcommand{\A}{\mathbb{A}}
\newcommand{\PP}{\mathbb{P}}
\newcommand{\E}{\mathbb{E}}
\newcommand{\eps}{\varepsilon}
\newcommand{\dd}{\,\mathrm{d}}
\begin{document}
\section*{UnsolvedMath ID: 30002991}
\textbf{Source problem code:} catalogue entry
\subsection{Definitions and mathematical statement}
The supplied question uses a base set \(B\), a target \(M\) with tangent bundle \(TM\), a domain dimension \(m\), and a multiplicity \(Q\). The notation implicitly requires positive integers \(m,Q\) and a differentiable structure on \(M\), but the record does not specify further hypotheses on these objects. Interpret the tangent-map bundle as
\[
\operatorname{Hom}(\R^m,TM)
=\bigsqcup_{\nu\in M}\operatorname{Hom}(\R^m,T_\nu M).
\]
Thus an argument of the integrand consists of \(x\in B\), \(\nu\in M\), and a linear map \(N:\R^m\to T_\nu M\).

At the level of sets, the space of unordered \(Q\)-tuples in \(M\), with repetitions allowed, is
\[
\mathcal A_Q(M)=M^Q/S_Q,
\]
where \(S_Q\) is the permutation group acting on the entries. A \(Q\)-valued map into \(M\) takes values in \(\mathcal A_Q(M)\). This set-theoretic definition does not specify a metric, a function space, or a notion of derivative for such maps.

\textbf{Question, with unresolved source qualifications.} Suppose that
\[
f:B\times\operatorname{Hom}(\R^m,TM)\longrightarrow[0,+\infty)
\]
is quasiconvex in the sense intended by the source. For \(x\in B\), \(\nu_i\in M\), and \(N_i\in\operatorname{Hom}(\R^m,T_{\nu_i}M)\), \(1\leq i\leq Q\), define
\[
e_M(x,\nu_1,\ldots,\nu_Q,N_1,\ldots,N_Q)
:=\sum_{i=1}^Q f(x,\nu_i,N_i).
\]
The pairs \((\nu_i,N_i)\) are considered up to simultaneous permutation. Must \(e_M\) be quasiconvex on the space of \(Q\)-valued maps into \(M\)?

The supplied record does not define either quasiconvexity predicate or the admissible variations, comparison domains, and boundary conditions. It also leaves the geometry of \(M\), the role and domain properties of \(B\), and any analytic assumptions on \(f\) unspecified. Consequently, this is a faithful conditional restatement of the supplied question; a fully quantified variational inequality cannot be fixed from the supplied evidence alone.
\subsection{Short English statement}
Does summing a nonnegative quasiconvex integrand over the \(Q\) entries of a multiple-valued map preserve quasiconvexity? The supplied record labels this question open, but its precise variational definitions and hypotheses are missing. This statement-only draft requires human mathematical review.
\subsection{Sources}
\begin{itemize}
\item Local UnsolvedMath registry, permanent ID \(30002991\), source problem code \texttt{catalogue entry}, title: \emph{Quasiconvexity of Multiple-Valued Integrands}. The supplied record belongs to the Oberwolfach Reports Open Problems collection. No individual proposer is supplied.
\item Supplied report citation: \emph{Mini-Workshop: Scales in Plasticity (2016)}, Oberwolfach Reports, pp.\ 3007--3048, DOI \texttt{10.4171/owr/2015/53}. \url{https://doi.org/10.4171/owr/2015/53}. The record gives proposed year 2015; the relationship between that year and the citation's 2016 designation is unresolved.
\item Supplied literature source: \emph{Oberwolfach Q-valued report}. \url{https://ems.press/content/serial-article-files/46600}.
\item Supplied literature source: \emph{Uniform polyconvexity and Q-graphs}. \url{https://arxiv.org/abs/2603.20788}.
\item Supplied literature source: \emph{Q-valued variational regularity}. \url{https://arxiv.org/abs/0803.0060}.
\item Supplied literature assessment dated \texttt{2026-08-21}: the extracted sentence is described as a fragment of the general \(Q\)-valued quasiconvexity question; special structured cases are reported as known, and the curation reports finding no general resolution or counterexample for the stated problem. The supplied status is \texttt{open}. These are supplied dated curation claims, not independently verified mathematical or literature conclusions.
\item The registry labels the statement \texttt{corrected\_verified} and reports that the official report passage was checked directly and its mathematical question, definitions, and notation restored without strengthening the claim. This is a supplied verification assertion; no source visit or independent verification is claimed here.
\end{itemize}
\end{document}
